% =====================================================================
%  Explicación sencilla de los gráficos del experimento de ergodicidad
%  Archivo independiente: compilar con pdflatex (dos pasadas)
% =====================================================================
\documentclass[11pt,a4paper]{article}

\usepackage[utf8]{inputenc}
\usepackage[T1]{fontenc}
\usepackage[spanish]{babel}
\usepackage{amsmath,amssymb}
\usepackage{geometry}
\usepackage{xcolor}
\usepackage{tcolorbox}
\usepackage{enumitem}
\usepackage{booktabs}
\usepackage{hyperref}

\geometry{margin=2.2cm}
\hypersetup{colorlinks=true,linkcolor=blue!50!black,urlcolor=blue!50!black}

\definecolor{cPi}{RGB}{214,69,65}      % π₀⁽¹⁾
\definecolor{cPi}{RGB}{232,163,47}     % π₀⁽²⁾
\definecolor{cPi}{RGB}{60,140,220}     % π₀⁽³⁾
\definecolor{idea}{RGB}{40,40,40}

\newcommand{\E}{|E(G_t)|}
\newcommand{\cond}[3]{\textcolor{#1}{\pi_0^{(#2)}}\,:#3}

\tcbset{
  myblock/.style={
    colback=white,
    colframe=blue!40!black,
    fonttitle=\bfseries,
    arc=2mm,
    boxrule=0.8pt,
    left=6pt,right=6pt,top=4pt,bottom=4pt
  }
}

\title{Guía de lectura de los gráficos del experimento de ergodicidad}
\author{Ergodicidad y convergencia en redes espacio-temporales}
\date{}

\begin{document}
\maketitle

% ---------------------------------------------------------------------
\begin{tcolorbox}[myblock,title=El experimento en tres frases]
Tres simulaciones corren \textbf{en paralelo} con el mismo modelo
\((X,\{C_i^X\},k,\alpha)\) pero con \textbf{tres puntos de partida distintos}:

\begin{center}
\begin{tabular}{lll}
\toprule
\textbf{Condición inicial} & \textbf{Grafo de partida} & \textbf{Color de partida} \\
\midrule
\cond{red}{1}{\(K_n+C_1\)}   & completo (todo conectado) & todos los vértices en \(C_1\) \\
\cond{orange}{2}{\(K_{1,n-1}+C_2\)} & estrella            & todos los vértices en \(C_2\) \\
\cond{blue}{3}{\(K_n+C_0\)}  & completo (todo conectado) & todos los vértices en \(C_0\) \\
\bottomrule
\end{tabular}
\end{center}

La pregunta es: \textbf{¿el punto de partida cambia el resultado al final?}
Si la respuesta es \emph{no}, decimos que el sistema es \textbf{ergódico}:
existe una única distribución de equilibrio \(\pi^*\) y se llega a ella
desde cualquier estado inicial \(\pi_0\).
\end{tcolorbox}

% ---------------------------------------------------------------------
\begin{tcolorbox}[myblock,title=1. Media de \(\E\) por condición inicial --- \textit{``¿hacia dónde camina cada simulación?''}]
\textbf{Qué muestra.} Una línea por condición inicial. El eje horizontal es el
tiempo (los pasos \(t\)) y el vertical es el \textbf{número promedio de aristas}
del grafo en ese instante, promediado entre las 10 cadenas de cada condición.

\textbf{Cómo leerlo.} Al inicio las tres líneas están separadas (la dispersión
es de \textbf{10 aristas}): las condiciones de partida son muy distintas
(\(K_n\) tiene \(\binom{n}{2}\) aristas, la estrella solo \(n-1\)). A medida que
avanza el tiempo, las tres líneas se \textbf{mezclan} y quedan ondulando juntas
en la misma banda, alrededor de \(\sim5\) aristas.

\textbf{Para qué sirve.} Es la prueba visual de que las tres carreras
\textbf{acaban en el mismo sitio}. El sistema olvida cómo empezó.
\end{tcolorbox}

% ---------------------------------------------------------------------
\begin{tcolorbox}[myblock,title=2--4. Frecuencia de \(C_0\), \(C_1\) y \(C_2\) --- \textit{``¿de qué color se pinta cada cosa?''}]
\textbf{Qué muestra.} Tres gráficos (uno por color). El eje vertical va de 0 a 1
y representa el \textbf{porcentaje de vértices} que tienen ese color en ese
instante. Por ejemplo, una línea en \(0{,}77\) significa que el \(77\%\) de los
puntos son de ese color.

\textbf{Cómo leerlo.} En cada gráfico se superponen las tres condiciones. Al
final todas se estabilizan en los mismos valores: un color domina
(\(\approx0{,}77\)) y los otros dos rondan \(\approx0{,}12\) cada uno.

\textbf{Para qué sirve.} Muestra que la \textbf{mezcla de colores} alcanza el
mismo equilibrio sin importar con qué color empezaste. Nótese que esto es un
equilibrio \emph{dinámico}: las líneas siguen temblando, pero su nivel promedio
no cambia.
\end{tcolorbox}

% ---------------------------------------------------------------------
\begin{tcolorbox}[myblock,title=5. Distribución de \(\E\) en la ventana final --- \textit{``la fotografía del final''}]
\textbf{Qué muestra.} Un \textbf{histograma}: las barras de la última ventana
(últimos \(750\) pasos de las \(3000\)). Cada barra responde: ``¿cuántas veces
apareció un grafo con \(m\) aristas?''. La altura total representa cuántas
veces se observó esa cantidad de aristas.

\textbf{Cómo leerlo.} La barra más alta está en \(5\) aristas: es lo más
frecuente. Las tres condiciones dan barras \textbf{casi idénticas}, superpuestas,
con media global de \(5{,}1\) aristas.

\textbf{Para qué sirve.} Aquí el argumento es más fuerte que en el gráfico 1:
no solo coinciden los \emph{promedios}, sino que coincide \textbf{toda la forma
de la distribución}. Es la evidencia de que se llegó a la misma \(\pi^*\).
\end{tcolorbox}

% ---------------------------------------------------------------------
\begin{tcolorbox}[myblock,title=6. Dispersión entre condiciones (\(\max-\min\) de \(\langle\E\rangle\)) --- \textit{``el cronómetro del olvido''}]
\textbf{Qué muestra.} Una sola línea gris. En cada instante se toman las tres
medias de aristas y se calcula
\[
\max_{i}\langle\E\rangle_i \;-\; \min_{i}\langle\E\rangle_i,
\]
es decir, \textbf{cuánto se separan la condición ``mejor'' y la ``peor''} en ese
momento.

\textbf{Cómo leerlo.} Al principio la línea dispara (hasta \(\approx4{,}5\)
aristas de diferencia) porque las condiciones iniciales son distintas. Después
se queda \textbf{cerca de cero}, con vaivenes chiquitos.

\textbf{Para qué sirve.} Es el \textbf{medidor numérico de la ergodicidad}:
\begin{itemize}[nosep]
  \item Dispersión al inicio: \(10\) aristas \(\Rightarrow\) las condiciones son
        genuinamente distintas.
  \item Dispersión en la ventana final: \(0{,}15\) aristas \(\Rightarrow\) casi
        idénticas, dentro del ruido estadístico.
\end{itemize}
Si la línea se queda pegada abajo, las tres condiciones ya dicen lo mismo y el
estado inicial fue olvidado.
\end{tcolorbox}

% ---------------------------------------------------------------------
\begin{tcolorbox}[myblock,title=Los seis gráficos dicen lo mismo, desde ángulos distintos]
\begin{center}
\begin{tabular}{p{0.28\linewidth}p{0.62\linewidth}}
\toprule
\textbf{Gráfico} & \textbf{Qué responde en una frase} \\
\midrule
1. Media de \(\E\) & ¿Hacia dónde camina la simulación? (todas, al mismo sitio) \\
2--4. Frecuencias de \(C_0,C_1,C_2\) & ¿Qué mezcla de colores queda al final? (la misma) \\
5. Distribución de \(\E\) & ¿Se ve igual toda la distribución final? (sí, no solo la media) \\
6. Dispersión \(\max-\min\) & ¿Cuánto tarda en olvidar el inicio? (se desploma a \(\approx0\)) \\
\bottomrule
\end{tabular}
\end{center}

\medskip
\textbf{Conclusión.} El equilibrio alcanzado es \textbf{estadístico, no
estático}: el diagrama nunca deja de cambiar, pero la frecuencia con la que
visita cada configuración converge a una distribución fija \(\pi^*\),
\emph{completamente independiente de la red y la coloración iniciales}:
\[
\lim_{t\to\infty}\pi_0 P^{t}=\pi^{*}.
\]
\end{tcolorbox}

\end{document}
